% ============================================================================
%  SECTION 3 — Kinematic growth
% ============================================================================
\begin{frame}[t]{}
  \vfill\begin{center}{\Huge\color{YaleBlue} 3.\ \ Kinematic growth}\\[0.5em]
  {\color{BoxGray} Rodriguez, Hoger \& McCulloch (1994)}\end{center}\vfill
\end{frame}

\begin{frame}[t]{The idea, and the kinematics}
  Postulate a stress-free \textbf{growth} deformation; only the elastic part
  carries load. Growth is one scalar $\theta$ (radial thickening):
  \giveneq{\bm{F} = \bm{F}_e\,\bm{F}_g,\qquad \frac{M}{M_0}=\theta,\qquad
    \text{only }\bm{F}_e\text{ stores energy.}}
  \bigfigw{kin_multiplicative_split.png}{0.62\linewidth}{One growth variable for
  the whole mixture --- no individual turnover. That is the strength and the
  ceiling.}
  \figcredit{Fig.: A. Gebauer (TUM lecture); Rodriguez, Hoger, McCulloch, J Biomech (1994)}
\end{frame}

\begin{frame}[t]{The growth law}
  Growth nulls the stress deviation from the set-point (first-order, production
  $\propto$ current mass):
  \slboxeq{2.6em}{eq:kingrowth}{\frac{\dd\theta}{\dd t}
    = k_g\,\theta\left(\frac{\bar\sigma}{\bar\sigma_h}-1\right)}
  The set-point $\bar\sigma_h$ is \textbf{prescribed}: wherever growth stops
  ($\dot\theta=0$), stress is homeostatic \emph{by construction}. At each instant
  $\lambda$ follows from equilibrium
  $\bar\sigma(\lambda)=\sigma_{\text{req}}(\lambda,\theta)$.
  \srccite{Rodriguez, Hoger, McCulloch, J Biomech (1994)}
\end{frame}

\begin{frame}[t]{Stability is imposed, not predicted}
  Growth is stable iff adding mass \emph{lowers} the stress it must carry,
  $\left.\dd\bar\sigma/\dd\theta\right|_{\theta^*}<0$ --- always true for the
  \textbf{bar}, but a strong enough insult flips the sign for the \textbf{artery}
  (Laplace $\lambda^2$) and growth runs away.
  \bigfigw{fig02_kinematic_growth.pdf}{0.78\linewidth}{Kinematic growth either
  converges to the \emph{prescribed} set-point or diverges --- it cannot
  \emph{predict} stability from turnover.}
  \figcredit{Reproduced with the \texttt{gr} package (this repo)}
\end{frame}
